\begin{abstract}
We study short relations in the quadratic algebra
$(\mathbb Z/N\mathbb Z)[X]/(X^2-b)$.  Suppose $b$ is a square modulo
one prime divisor $p$ of the odd integer $N$, a nonsquare modulo another
prime divisor $q$, and an even supplied exponent $M$ is divisible by
$p-1$.  A coefficient-gcd chain finds a proper factor with probability
at least $(p+3)/(2p)>1/2$, using $O(\log M)$ quadratic-algebra operations.
Other prime factors and arbitrary prime-power multiplicities are
allowed.  The proof continues to work when $M$ annihilates both local
quadratic unit groups, a case in which the one-shot relation can be
uninformative.  For a distinct-prime semiprime, the resulting factor
also constructs the actual selected value and a logarithmic-size
division-free section circuit of an earlier succinct rank-three bundle.
We prove this transfer directly from its chart equations.  The result
is an auxiliary supplied-input theorem, not a new general-factorization
exponent.
\end{abstract}

\tableofcontents

\section{A short relation with a precise input promise}

A relation in a small algebra is useful for factoring when its
coefficients vanish in one hidden prime component but not another.
A short expression alone does not guarantee this discrepancy.
The theorem below gives both an explicit expression and a proof of
uniform success under a supplied-exponent promise.

The core algorithm and its proof are in Sections~\ref{sec:qr-algorithm}
and~\ref{sec:qr-success}.  Section~\ref{sec:qr-saturation} explains why
the intermediate chain matters.  Section~\ref{sec:qr-section} gives the
selected-section transfer, including its explicit circuit.

\begin{theorem}[Mixed-prime quadratic relation chain]
\label{thm:qr-main}
Let $N$ be odd, let $b$ be a unit modulo $N$, and suppose there are
distinct primes $p,q\mid N$ with
\[
 \left(\frac bp\right)=1,\qquad
 \left(\frac bq\right)=-1.
\]
Supply a positive even integer $M$ satisfying $p-1\mid M$.
The primes need not be known to the algorithm.  There is a randomized
algorithm which, on each independent trial, returns a verified proper
factor of $N$ with probability at least
\[
 \frac{p+3}{2p}>\frac12,
\]
and otherwise reports failure.  A trial uses $O(\log M)$ quadratic-ring
multiplications and gcd operations, hence
$O(\operatorname{bitlen}(M)(\log N)^2)$ elementary bit operations.
The theorem permits repeated prime factors and additional prime divisors
of $N$.  It does not require $q-1\mid M$ or a factorization of $M$.
\end{theorem}

Repeated trials give a Las Vegas algorithm with fewer than two expected
trials: every returned factor is checked, and only the running time is
random.  The assumption on $M$ is a genuine supplied input.  If it is
constructed by another algorithm, that construction and its output
bitlength must be charged separately.

This develops the earlier short-root-algebra relation specification
N1101 and its selected-section specialization N1505 in the retained
catalogue \cite{BroxCatalogue2026}.  Randomized factor extraction from a
supplied order multiple is classical; the contribution here is the
explicit quadratic coefficient construction, its saturation-resistant
proof, and the checked transfer to that particular section target.
No claim of a new randomized general-factoring complexity class is made.

\section{The algorithm uses two coefficients}
\label{sec:qr-algorithm}

Set
\[
 R=\mathbb Z/N\mathbb Z,
 \qquad A=R[X]/(X^2-b).
\]
An element is represented exactly as $c+dX$.  Multiplication is
\begin{equation}
 (c+dX)(c'+d'X)=(cc'+bdd')+(cd'+dc')X,
 \label{eq:qr-multiply}
\end{equation}
so its dimension is fixed and its coordinates have $O(\log N)$ bits.
Its conjugate is $c-dX$ and its norm is $c^2-bd^2$.

\begin{algorithm}[One coefficient-chain trial]
\label{alg:qr-chain}
Input $N,b,M$ as in Theorem~\ref{thm:qr-main}.
\begin{enumerate}
\item Sample $a$ uniformly from $\{0,\ldots,N-1\}$.
      Compute $g=\gcd(a^2-b,N)$ and return it if $1<g<N$.
\item Write $M=2^tu$ with $u$ odd, and compute
      $z_0=(X-a)^u\pmod{X^2-b,N}$ by binary powering.
\item For $j=0,\ldots,t$, write $z_j=c_j+d_jX$ and test
      \[
       \gcd(N,d_j),\qquad \gcd(N,c_j-1,d_j).
      \]
      Return any proper factor found.  If $j<t$, set
      $z_{j+1}=z_j^2$ using \eqref{eq:qr-multiply}.
\item If no proper gcd was found, report failure for this trial.
\end{enumerate}
\end{algorithm}

Under the stated promise, $\gcd(a^2-b,N)$ cannot equal $N$: the
nonsquare condition modulo $q$ makes $a^2-b$ nonzero modulo $q$ for
every $a$.  An implementation may nevertheless treat this value as an
unsuccessful trial on inputs outside the promise.
Exact uniform sampling takes an expected bounded number of
$\operatorname{bitlen}(N)$-bit draws by rejection.
The initial powering and the entire squaring chain together use
$O(\operatorname{bitlen}(M))$ ring operations.  The two gcd tests per
stage use the same bit bound.  Nothing of degree $M$ is expanded.

\section{Why more than half the samples succeed}
\label{sec:qr-success}

\begin{proof}[Proof of Theorem~\ref{thm:qr-main}]
Choose $r\in\mathbb F_p^*$ with $r^2=b$.  The quadratic algebra
modulo $p$ is the product of two fields, with
\[
 X-a\longmapsto(r-a,-r-a).
\]
Modulo $q$ it is the field $\mathbb F_{q^2}$.
Put $e=v_2(p-1)$, and write $p-1=2^em$ with $m$ odd.
Because $p-1\mid M=2^tu$, one has $e\leq t$ and $u/m$ odd.

Suppose first
\begin{equation}
 \left(\frac{a^2-b}{p}\right)=-1.
 \label{eq:qr-good-a}
\end{equation}
The two nonzero split coordinates have opposite quadratic characters,
since their product is $a^2-b$.
At stage $e-1$, raising to $u2^{e-1}$ therefore gives the two opposite
values $1,-1$; Euler's criterion applies because $u/m$ is odd.
At stage $e$ both coordinates are one.  Consequently
\begin{equation}
 d_{e-1}\ne0\pmod p,
 \qquad c_e-1=d_e=0\pmod p.
 \label{eq:qr-two-levels}
\end{equation}

In the field modulo $q$ there are two possibilities.
If $z_e\ne1$, at least one of $c_e-1,d_e$ is nonzero modulo $q$.
The content gcd $\gcd(N,c_e-1,d_e)$ contains $p$ and excludes $q$,
so it is proper.  If $z_e=1$, then $z_{e-1}^2=1$ in a field of odd
characteristic, so $z_{e-1}=1$ or $-1$.  Its linear coefficient is zero
modulo $q$ but nonzero modulo $p$ by \eqref{eq:qr-two-levels}.
Thus $\gcd(N,d_{e-1})$ contains $q$ and excludes $p$, again proper.
Other primes or larger prime-power multiplicities cannot invalidate
either conclusion: a gcd containing one prime and excluding another
is neither one nor $N$.

The two residues $a=\pm r\pmod p$ also always succeed, through
the first norm gcd, because $a^2-b$ is never zero modulo $q$.
It remains to count \eqref{eq:qr-good-a}.  Counting pairs $(a,y)$ with
$y^2=a^2-b$ gives $p-1$ pairs: the equivalent equation
$(a-y)(a+y)=b$ has one pair for each nonzero first factor, since 2 is
invertible.  Therefore
\[
 \sum_{a\bmod p}\left(\frac{a^2-b}{p}\right)=-1.
\]
There are two zeros, hence exactly $(p-1)/2$ negative values.
Including the two norm-zero residues gives $(p+3)/2$ guaranteed
successful residues modulo $p$.  A uniform residue modulo $N$
projects uniformly to $\mathbb F_p$, proving the claimed probability.
The cost follows from Algorithm~\ref{alg:qr-chain}.
\end{proof}

\begin{corollary}[The distinct-prime semiprime promise]
\label{cor:qr-semiprime}
Let $N=pq$ be a product of distinct odd primes and supply an even $M$
divisible by both $p-1$ and $q-1$.  A unit $b$ with Jacobi symbol
$(b/N)=-1$ satisfies the mixed-type hypothesis after relabeling the
primes.  Such a $b$ can itself be found by random sampling and Jacobi
tests in expected polynomial time in $\log N$, or a norm gcd may
already produce a factor.  Thus the complete supplied-exponent
quadratic-relation algorithm has expected polynomial bit cost in
$\log N+\operatorname{bitlen}(M)$.
\end{corollary}

\begin{proof}
The two Legendre signs have product $-1$, so one is positive and one
negative.  Exactly half of the units modulo a distinct-prime semiprime
have Jacobi sign $-1$, by the Chinese remainder theorem and the two
equal-sized square/nonsquare classes in either field.
Their density among all residues is
$\frac12(1-1/p)(1-1/q)\geq4/15$, since distinct odd primes are at
least 3 and 5.  Gcd and Jacobi tests require polynomial bit cost and
do not require the factorization.  Apply Theorem~\ref{thm:qr-main}.
\end{proof}

For a general odd integer, Jacobi sign $-1$ alone does not guarantee
both a split and a nonsplit prime component.  The mixed-prime statement
of Theorem~\ref{thm:qr-main}, rather than that semiprime shortcut, is
the appropriate general promise.

\section{Why a saturated endpoint is not enough}
\label{sec:qr-saturation}

Consider the simpler test which uses only
$(X-a)^M-1$, together with its initial norm gcd.
For this section let $N=pq$ have the mixed types, and assume additionally
$q-1\mid M$.  In the nonsplit field choose $\alpha^2=b$; then
$\alpha^q=-\alpha$.  The map
\begin{equation}
 a\longmapsto v(a)=\frac{-\alpha-a}{\alpha-a}
                  =(\alpha-a)^{q-1}
 \label{eq:qr-norm-one-map}
\end{equation}
is a bijection from $\mathbb F_q$ to the norm-one group of
$\mathbb F_{q^2}^*$ with the identity removed.
Its values have norm one and are not one; injectivity follows by cross
multiplication, and the target has $q$ elements.
Equivalently its inverse is
$a=\alpha(v+1)/(v-1)$, fixed by Frobenius for norm-one $v$.

It follows that the number of $a\bmod q$ with $(\alpha-a)^M=1$ is
\begin{equation}
 d-1,\qquad d=\gcd(M/(q-1),q+1).
 \label{eq:qr-one-shot-zeros}
\end{equation}
Indeed this is the number of solutions to $v^{M/(q-1)}=1$ in a cyclic
group of order $q+1$, with its identity removed.  Thus the exact number
of successful residues $a\bmod N$ for the one-shot test is
\begin{equation}
 N-(p-2)(d-1).
 \label{eq:qr-one-shot-count}
\end{equation}
The failure choices are precisely the $p-2$ nonzero-norm residues at
the split prime and the $d-1$ vanishing-relation residues at the
nonsplit prime.

If $M$ is divisible by both $p^2-1$ and $q^2-1$, the relation vanishes
on every unit input, and only the initial norm check remains useful.
The chain in Algorithm~\ref{alg:qr-chain} still succeeds with
probability above one half.  It inspects the adjacent stages at which
the split coordinates acquire different signs, rather than relying
only on the final value.

As a finite exact illustration, for $N=35$, $b=2$, and $M=48$, the
split prime is 7 and the nonsplit prime is 5.
Formula \eqref{eq:qr-one-shot-count} gives exactly 10 successful
one-shot residues out of 35.  The supplied verification program
enumerates all 35 residues and finds a proper factor on every full-chain
trial.  That finite statement is an example, not the universal
probability proof, which was given above.

\section{The actual selected section of the earlier bundle}
\label{sec:qr-section}

We now connect the relation theorem to a specific earlier construction,
not merely to the general phrase ``selected scalar.''
Let $N=p_0q_0$ with distinct odd primes $p_0<q_0$, let
$R=\mathbb Z/N\mathbb Z$, and set
\[
 t_0=\lfloor\sqrt N\rfloor,
 \qquad m=\lfloor t_0/2\rfloor+1,
 \qquad G(X)=(1+X)^N-1.
\]
On the two standard charts of the projective line use the rank-two
transition
\begin{equation}
 T_N(X)=\begin{pmatrix}X^{-m}&X^{-m}G(X)\\0&X^m\end{pmatrix}.
 \label{eq:qr-bundle}
\end{equation}
Adding the line bundle of degree $m$ gives the earlier rank-three
construction.  Its selected coordinate is the second coordinate at
infinity of a section of the rank-two bundle twisted by $-m$.
The Laurent transition has an $O(\log N)$-gate expression by binary
powering; this succinctness by itself does not construct a section.

The selected coordinate and its ideal can be determined directly from
the chart equations, without a splitting theorem for bundles.
Write a section on the finite chart as $(u_0(X),c_0(X))$ and on the
infinity chart as $(u_\infty(X^{-1}),c_\infty(X^{-1}))$.
The twisted transition equation is
\[
 \binom{u_0}{c_0}
 =\begin{pmatrix}X^{-2m}&X^{-2m}G(X)\\0&1\end{pmatrix}
   \binom{u_\infty}{c_\infty}.
\]
The second coordinate must lie in
$R[X]\cap R[X^{-1}]=R$, so it is a constant $c$.
The first equation becomes
\begin{equation}
 X^{2m}u_0=u_\infty+cG(X).
 \label{eq:qr-gluing}
\end{equation}
Comparing nonpositive powers forces $u_\infty=0$, since $G$ has zero
constant term.  Comparing powers $1,\ldots,2m-1$ shows that a section
exists exactly when
\begin{equation}
 c\binom Nj=0\pmod N\qquad(1\leq j\leq2m-1).
 \label{eq:qr-prefix-annihilator}
\end{equation}

\begin{proposition}[The selected ideal is exactly $p_0R$]
\label{prop:qr-selected}
The possible constants $c$ in \eqref{eq:qr-gluing} form the ideal
$p_0R$.  Every nonzero section has a nonzero such constant, and
$\gcd(c,N)=p_0$ for every nonzero $c\in p_0R$.
\end{proposition}

\begin{proof}
We have
$p_0\leq t_0\leq2m-1\leq t_0+1<q_0$.
The last inequality follows from
$p_0\leq q_0-2$ and $p_0q_0<(q_0-1)^2$.
Modulo $q_0$, Frobenius gives
$(1+X)^N=(1+X^{q_0})^{p_0}$, so every coefficient in the prefix
$1\leq j\leq2m-1$ vanishes.
Modulo $p_0$, the coefficient at $j=p_0$ equals $q_0\ne0$.
The prefix coefficient ideal in $R$ is therefore $q_0R$, and its
annihilator is $p_0R$.  Equation \eqref{eq:qr-prefix-annihilator}
proves the first claim.  If $c=0$, \eqref{eq:qr-gluing} forces
$u_0=u_\infty=0$, so the section is zero.  A nonzero element of
$p_0R$ is divisible by $p_0$ and not by $q_0$, proving the gcd claim.
\end{proof}

\begin{proposition}[A division-free section from one selected value]
\label{prop:qr-section-circuit}
Given any nonzero $c\in p_0R$, one can construct a complete section
circuit with $O(\log N)$ arithmetic gates.  Conversely, a nonzero
section's selected value gives a proper factor in polynomial time.
\end{proposition}

\begin{proof}
Compute $p_0=\gcd(c,N)$ and $q_0=N/p_0$, and write
$S_k(U)=1+U+\cdots+U^{k-1}$.
The chart vectors are
\begin{equation}
 s_\infty=\binom0c,
 \qquad
 s_0=\binom{cX^{q_0-2m}S_{p_0}(1+X^{q_0})}{c}.
 \label{eq:qr-section-circuit}
\end{equation}
The exponent is nonnegative because $2m-1<q_0$.
Modulo $p_0$ all relevant multiples of $c$ vanish.
Modulo $q_0$,
\[
 cG(X)=c((1+X^{q_0})^{p_0}-1)
       =cX^{q_0}S_{p_0}(1+X^{q_0}).
\]
Chinese remaindering proves \eqref{eq:qr-gluing} over $R$.
Joint binary recursion computes $U^k$ and $S_k(U)$ using
\[
 S_{2k}(U)=S_k(U)(1+U^k),\qquad
 S_{2k+1}(U)=S_{2k}(U)+U^{2k}.
\]
Together with binary powering for the displayed monomials, this gives
the gate bound and uses no polynomial division.
The converse is Proposition~\ref{prop:qr-selected}.
\end{proof}

Theorem~\ref{thm:qr-main} supplies a factor $h$ on its promised
semiprime input.  Set $p_0=\min(h,N/h)$ and $c=p_0\pmod N$;
Proposition~\ref{prop:qr-section-circuit} then constructs the
\emph{actual requested section}.
This explicitly reuses the previous selected-section reconstruction,
which has just been reproved from its equations.
It does not give a fast section constructor from
\eqref{eq:qr-bundle} alone: the supplied exponent and the useful local
type are extra data.  Factoring equivalence is a description of this
interface, not a reason to reject other possible section constructors.

\section{Two additional exact readouts and their limits}

\subsection{A local rank witness from two characteristic coefficients}

Multiplication by $z_j=c_j+d_jX$ is the matrix
\[
 U_j=\begin{pmatrix}c_j&bd_j\\d_j&c_j\end{pmatrix},
 \quad \tau_j=2c_j,\quad\delta_j=c_j^2-bd_j^2.
\]
The two quantities
\[
 D_j=1-\tau_j+\delta_j=\det(U_j-I),
 \qquad E_j=\tau_j^2-4\delta_j=4bd_j^2
\]
give the same success guarantee by testing their gcds with $N$.
At the later decisive stage in \eqref{eq:qr-two-levels},
$U_j-I$ is zero modulo $p$; if the element is not one in the nonsplit
field modulo $q$, that multiplication operator is invertible there,
so $D_j$ separates.  Otherwise $E_{j-1}$ separates because
$d_{j-1}$ vanishes only at the nonsplit one of these two primes.
This is an explicit constant-dimensional rank witness, not a generic
high-dimensional implicit-rank algorithm.

For comparison, over $\mathbb F_2$ the symmetric matrix
$\bigl(\begin{smallmatrix}1&1\\1&1\end{smallmatrix}\bigr)$ has rank
one and square zero.  It and the zero matrix have the same characteristic
polynomial $T^2$ and every positive power trace zero.  Thus generic
kernel dimensions are not determined by those invariants alone.
This does not exclude specially structured matrix algorithms.

\subsection{A quotient coefficient is not an arbitrary prefix}

Binary powering computes
\[
 (X-a)^E\bmod(X^2-b)=c_E+d_EX,
 \quad
 d_E=\sum_{k\text{ odd}}\binom Ek(-a)^{E-k}b^{(k-1)/2},
\]
in $O(\log E)$ ring operations.  The squaring recurrences are
$c_{2E}=c_E^2+bd_E^2$ and $d_{2E}=2c_Ed_E$.
This is an inexpensive selected coefficient in a fixed-degree quotient.
It does not recover a general truncation: $X^h(X^2-b)$ has quotient
value zero but prefix through degree $h$ equal to $-bX^h$, with unit
coefficient.  The local exponent theorem, not the quotient operation
alone, makes the coefficients here factor-sensitive.

\section{Reproducibility and unresolved construction costs}

The implementation takes only $N,b,M,a$ and uses exact modular
arithmetic and gcds.  Known prime factors occur in the verification
harness solely to check the universal local success implication on
finite examples, including $p^2q$, $pq^2$, and additional prime
components.  The one-shot formula \eqref{eq:qr-one-shot-count}, the
matrix identities, and the section gluing equations have separate exact
checks.  Such finite tests supplement, and do not replace, the proofs.

The unresolved costs are specific.  General factoring still needs an
appropriate supplied exponent or a different useful relation.  The
semiprime mixed type can be obtained by randomized Jacobi sampling, but
a deterministic sufficiently short successful list is not supplied by
this coefficient-chain argument.  The selected section is constructed
under that same extra promise, not from its succinct transition alone.
No claim of a general algorithm with first $L$-parameter below $1/3$
is made.
